% Môn: Vật lý
% Lớp: 10
% Chương: Chương II
% Bài: L10C2 Bài 7. Đồ thị dịch chuyển thời gian
% Loại: Lý thuyết
% File riêng, không trộn vào de.tex
% Định nghĩa lệnh Lý thuyết — dùng khi biên dịch lt.tex bằng TeX.
% App web đọc lt.tex trực tiếp (bỏ \input file này).
% Trong lt.tex, từ thư mục bài:
%   % Định nghĩa lệnh Lý thuyết — dùng khi biên dịch lt.tex bằng TeX.
% App web đọc lt.tex trực tiếp (bỏ \input file này).
% Trong lt.tex, từ thư mục bài:
%   % Định nghĩa lệnh Lý thuyết — dùng khi biên dịch lt.tex bằng TeX.
% App web đọc lt.tex trực tiếp (bỏ \input file này).
% Trong lt.tex, từ thư mục bài:
%   \input{../../../../_lenh/lythuyet.tex}
% từ gốc repo:
%   \input{ngan-hang/_lenh/lythuyet.tex}
\makeatletter
\providecommand{\indam}[1]{\textbf{#1}}
\providecommand{\chuy}[1]{\par\smallskip\noindent\textit{#1}\par}
\providecommand{\luuy}[1]{\par\smallskip\noindent\textbf{Lưu ý.} #1\par}
\providecommand{\ghichu}[1]{\par\smallskip\noindent\textbf{Ghi chú.} #1\par}
\providecommand{\vidu}[1]{\par\smallskip\noindent\textbf{Ví dụ.} #1\par}
\providecommand{\Blythuyet}{\subsection*{Lý thuyết}}
% Hai cột: kiến thức | hình
\providecommand{\haicotchay}[2]{%
  \par\medskip\noindent
  \begin{minipage}[t]{0.52\linewidth}#1\end{minipage}\hfill
  \begin{minipage}[t]{0.46\linewidth}#2\end{minipage}%
  \par\medskip
}
\providecommand{\kienthuccot}[2]{\haicotchay{#1}{#2}}
% Dự phòng nếu chưa nạp graphicx / caption (không ghi đè bản thật)
\providecommand{\resizebox}[3]{#3}
\providecommand{\captionof}[2]{\par\smallskip\noindent\textit{#2}\par}
\newenvironment{hoatdong}[1]{%
  \par\medskip\noindent\textbf{Hoạt động.} #1\par\smallskip
}{\par\medskip}
\newenvironment{emcobiet}{%
  \par\medskip\noindent\textbf{Em có biết}\par\smallskip
}{\par\medskip}
\newenvironment{emdahoc}{%
  \par\medskip\noindent\textbf{Em đã học}\par\smallskip
}{\par\medskip}
\newenvironment{emcothe}{%
  \par\medskip\noindent\textbf{Em có thể}\par\smallskip
}{\par\medskip}
\newenvironment{kienthuc}{%
  \par\medskip\noindent\textbf{Kiến thức cốt lõi}\par\smallskip
}{\par\medskip}
\newenvironment{traloi}{%
  \par\smallskip\noindent\textit{Trả lời / ghi chú của em}\par
  \vspace{1.2em}%
}{\par\medskip}
\newenvironment{phuongphap}{%
  \par\medskip\noindent\textbf{Phương pháp giải}\par\smallskip
}{\par\medskip}
\newenvironment{vidumau}{%
  \par\medskip\noindent\textbf{Bài mẫu}\par\smallskip
}{\par\medskip}
\newenvironment{dangmau}[1]{%
  \par\medskip\noindent\textbf{Dạng mẫu.} #1\par\smallskip
}{\par\medskip}
\makeatother

% từ gốc repo:
%   % Định nghĩa lệnh Lý thuyết — dùng khi biên dịch lt.tex bằng TeX.
% App web đọc lt.tex trực tiếp (bỏ \input file này).
% Trong lt.tex, từ thư mục bài:
%   \input{../../../../_lenh/lythuyet.tex}
% từ gốc repo:
%   \input{ngan-hang/_lenh/lythuyet.tex}
\makeatletter
\providecommand{\indam}[1]{\textbf{#1}}
\providecommand{\chuy}[1]{\par\smallskip\noindent\textit{#1}\par}
\providecommand{\luuy}[1]{\par\smallskip\noindent\textbf{Lưu ý.} #1\par}
\providecommand{\ghichu}[1]{\par\smallskip\noindent\textbf{Ghi chú.} #1\par}
\providecommand{\vidu}[1]{\par\smallskip\noindent\textbf{Ví dụ.} #1\par}
\providecommand{\Blythuyet}{\subsection*{Lý thuyết}}
% Hai cột: kiến thức | hình
\providecommand{\haicotchay}[2]{%
  \par\medskip\noindent
  \begin{minipage}[t]{0.52\linewidth}#1\end{minipage}\hfill
  \begin{minipage}[t]{0.46\linewidth}#2\end{minipage}%
  \par\medskip
}
\providecommand{\kienthuccot}[2]{\haicotchay{#1}{#2}}
% Dự phòng nếu chưa nạp graphicx / caption (không ghi đè bản thật)
\providecommand{\resizebox}[3]{#3}
\providecommand{\captionof}[2]{\par\smallskip\noindent\textit{#2}\par}
\newenvironment{hoatdong}[1]{%
  \par\medskip\noindent\textbf{Hoạt động.} #1\par\smallskip
}{\par\medskip}
\newenvironment{emcobiet}{%
  \par\medskip\noindent\textbf{Em có biết}\par\smallskip
}{\par\medskip}
\newenvironment{emdahoc}{%
  \par\medskip\noindent\textbf{Em đã học}\par\smallskip
}{\par\medskip}
\newenvironment{emcothe}{%
  \par\medskip\noindent\textbf{Em có thể}\par\smallskip
}{\par\medskip}
\newenvironment{kienthuc}{%
  \par\medskip\noindent\textbf{Kiến thức cốt lõi}\par\smallskip
}{\par\medskip}
\newenvironment{traloi}{%
  \par\smallskip\noindent\textit{Trả lời / ghi chú của em}\par
  \vspace{1.2em}%
}{\par\medskip}
\newenvironment{phuongphap}{%
  \par\medskip\noindent\textbf{Phương pháp giải}\par\smallskip
}{\par\medskip}
\newenvironment{vidumau}{%
  \par\medskip\noindent\textbf{Bài mẫu}\par\smallskip
}{\par\medskip}
\newenvironment{dangmau}[1]{%
  \par\medskip\noindent\textbf{Dạng mẫu.} #1\par\smallskip
}{\par\medskip}
\makeatother

\makeatletter
\providecommand{\indam}[1]{\textbf{#1}}
\providecommand{\chuy}[1]{\par\smallskip\noindent\textit{#1}\par}
\providecommand{\luuy}[1]{\par\smallskip\noindent\textbf{Lưu ý.} #1\par}
\providecommand{\ghichu}[1]{\par\smallskip\noindent\textbf{Ghi chú.} #1\par}
\providecommand{\vidu}[1]{\par\smallskip\noindent\textbf{Ví dụ.} #1\par}
\providecommand{\Blythuyet}{\subsection*{Lý thuyết}}
% Hai cột: kiến thức | hình
\providecommand{\haicotchay}[2]{%
  \par\medskip\noindent
  \begin{minipage}[t]{0.52\linewidth}#1\end{minipage}\hfill
  \begin{minipage}[t]{0.46\linewidth}#2\end{minipage}%
  \par\medskip
}
\providecommand{\kienthuccot}[2]{\haicotchay{#1}{#2}}
% Dự phòng nếu chưa nạp graphicx / caption (không ghi đè bản thật)
\providecommand{\resizebox}[3]{#3}
\providecommand{\captionof}[2]{\par\smallskip\noindent\textit{#2}\par}
\newenvironment{hoatdong}[1]{%
  \par\medskip\noindent\textbf{Hoạt động.} #1\par\smallskip
}{\par\medskip}
\newenvironment{emcobiet}{%
  \par\medskip\noindent\textbf{Em có biết}\par\smallskip
}{\par\medskip}
\newenvironment{emdahoc}{%
  \par\medskip\noindent\textbf{Em đã học}\par\smallskip
}{\par\medskip}
\newenvironment{emcothe}{%
  \par\medskip\noindent\textbf{Em có thể}\par\smallskip
}{\par\medskip}
\newenvironment{kienthuc}{%
  \par\medskip\noindent\textbf{Kiến thức cốt lõi}\par\smallskip
}{\par\medskip}
\newenvironment{traloi}{%
  \par\smallskip\noindent\textit{Trả lời / ghi chú của em}\par
  \vspace{1.2em}%
}{\par\medskip}
\newenvironment{phuongphap}{%
  \par\medskip\noindent\textbf{Phương pháp giải}\par\smallskip
}{\par\medskip}
\newenvironment{vidumau}{%
  \par\medskip\noindent\textbf{Bài mẫu}\par\smallskip
}{\par\medskip}
\newenvironment{dangmau}[1]{%
  \par\medskip\noindent\textbf{Dạng mẫu.} #1\par\smallskip
}{\par\medskip}
\makeatother

% từ gốc repo:
%   % Định nghĩa lệnh Lý thuyết — dùng khi biên dịch lt.tex bằng TeX.
% App web đọc lt.tex trực tiếp (bỏ \input file này).
% Trong lt.tex, từ thư mục bài:
%   % Định nghĩa lệnh Lý thuyết — dùng khi biên dịch lt.tex bằng TeX.
% App web đọc lt.tex trực tiếp (bỏ \input file này).
% Trong lt.tex, từ thư mục bài:
%   \input{../../../../_lenh/lythuyet.tex}
% từ gốc repo:
%   \input{ngan-hang/_lenh/lythuyet.tex}
\makeatletter
\providecommand{\indam}[1]{\textbf{#1}}
\providecommand{\chuy}[1]{\par\smallskip\noindent\textit{#1}\par}
\providecommand{\luuy}[1]{\par\smallskip\noindent\textbf{Lưu ý.} #1\par}
\providecommand{\ghichu}[1]{\par\smallskip\noindent\textbf{Ghi chú.} #1\par}
\providecommand{\vidu}[1]{\par\smallskip\noindent\textbf{Ví dụ.} #1\par}
\providecommand{\Blythuyet}{\subsection*{Lý thuyết}}
% Hai cột: kiến thức | hình
\providecommand{\haicotchay}[2]{%
  \par\medskip\noindent
  \begin{minipage}[t]{0.52\linewidth}#1\end{minipage}\hfill
  \begin{minipage}[t]{0.46\linewidth}#2\end{minipage}%
  \par\medskip
}
\providecommand{\kienthuccot}[2]{\haicotchay{#1}{#2}}
% Dự phòng nếu chưa nạp graphicx / caption (không ghi đè bản thật)
\providecommand{\resizebox}[3]{#3}
\providecommand{\captionof}[2]{\par\smallskip\noindent\textit{#2}\par}
\newenvironment{hoatdong}[1]{%
  \par\medskip\noindent\textbf{Hoạt động.} #1\par\smallskip
}{\par\medskip}
\newenvironment{emcobiet}{%
  \par\medskip\noindent\textbf{Em có biết}\par\smallskip
}{\par\medskip}
\newenvironment{emdahoc}{%
  \par\medskip\noindent\textbf{Em đã học}\par\smallskip
}{\par\medskip}
\newenvironment{emcothe}{%
  \par\medskip\noindent\textbf{Em có thể}\par\smallskip
}{\par\medskip}
\newenvironment{kienthuc}{%
  \par\medskip\noindent\textbf{Kiến thức cốt lõi}\par\smallskip
}{\par\medskip}
\newenvironment{traloi}{%
  \par\smallskip\noindent\textit{Trả lời / ghi chú của em}\par
  \vspace{1.2em}%
}{\par\medskip}
\newenvironment{phuongphap}{%
  \par\medskip\noindent\textbf{Phương pháp giải}\par\smallskip
}{\par\medskip}
\newenvironment{vidumau}{%
  \par\medskip\noindent\textbf{Bài mẫu}\par\smallskip
}{\par\medskip}
\newenvironment{dangmau}[1]{%
  \par\medskip\noindent\textbf{Dạng mẫu.} #1\par\smallskip
}{\par\medskip}
\makeatother

% từ gốc repo:
%   % Định nghĩa lệnh Lý thuyết — dùng khi biên dịch lt.tex bằng TeX.
% App web đọc lt.tex trực tiếp (bỏ \input file này).
% Trong lt.tex, từ thư mục bài:
%   \input{../../../../_lenh/lythuyet.tex}
% từ gốc repo:
%   \input{ngan-hang/_lenh/lythuyet.tex}
\makeatletter
\providecommand{\indam}[1]{\textbf{#1}}
\providecommand{\chuy}[1]{\par\smallskip\noindent\textit{#1}\par}
\providecommand{\luuy}[1]{\par\smallskip\noindent\textbf{Lưu ý.} #1\par}
\providecommand{\ghichu}[1]{\par\smallskip\noindent\textbf{Ghi chú.} #1\par}
\providecommand{\vidu}[1]{\par\smallskip\noindent\textbf{Ví dụ.} #1\par}
\providecommand{\Blythuyet}{\subsection*{Lý thuyết}}
% Hai cột: kiến thức | hình
\providecommand{\haicotchay}[2]{%
  \par\medskip\noindent
  \begin{minipage}[t]{0.52\linewidth}#1\end{minipage}\hfill
  \begin{minipage}[t]{0.46\linewidth}#2\end{minipage}%
  \par\medskip
}
\providecommand{\kienthuccot}[2]{\haicotchay{#1}{#2}}
% Dự phòng nếu chưa nạp graphicx / caption (không ghi đè bản thật)
\providecommand{\resizebox}[3]{#3}
\providecommand{\captionof}[2]{\par\smallskip\noindent\textit{#2}\par}
\newenvironment{hoatdong}[1]{%
  \par\medskip\noindent\textbf{Hoạt động.} #1\par\smallskip
}{\par\medskip}
\newenvironment{emcobiet}{%
  \par\medskip\noindent\textbf{Em có biết}\par\smallskip
}{\par\medskip}
\newenvironment{emdahoc}{%
  \par\medskip\noindent\textbf{Em đã học}\par\smallskip
}{\par\medskip}
\newenvironment{emcothe}{%
  \par\medskip\noindent\textbf{Em có thể}\par\smallskip
}{\par\medskip}
\newenvironment{kienthuc}{%
  \par\medskip\noindent\textbf{Kiến thức cốt lõi}\par\smallskip
}{\par\medskip}
\newenvironment{traloi}{%
  \par\smallskip\noindent\textit{Trả lời / ghi chú của em}\par
  \vspace{1.2em}%
}{\par\medskip}
\newenvironment{phuongphap}{%
  \par\medskip\noindent\textbf{Phương pháp giải}\par\smallskip
}{\par\medskip}
\newenvironment{vidumau}{%
  \par\medskip\noindent\textbf{Bài mẫu}\par\smallskip
}{\par\medskip}
\newenvironment{dangmau}[1]{%
  \par\medskip\noindent\textbf{Dạng mẫu.} #1\par\smallskip
}{\par\medskip}
\makeatother

\makeatletter
\providecommand{\indam}[1]{\textbf{#1}}
\providecommand{\chuy}[1]{\par\smallskip\noindent\textit{#1}\par}
\providecommand{\luuy}[1]{\par\smallskip\noindent\textbf{Lưu ý.} #1\par}
\providecommand{\ghichu}[1]{\par\smallskip\noindent\textbf{Ghi chú.} #1\par}
\providecommand{\vidu}[1]{\par\smallskip\noindent\textbf{Ví dụ.} #1\par}
\providecommand{\Blythuyet}{\subsection*{Lý thuyết}}
% Hai cột: kiến thức | hình
\providecommand{\haicotchay}[2]{%
  \par\medskip\noindent
  \begin{minipage}[t]{0.52\linewidth}#1\end{minipage}\hfill
  \begin{minipage}[t]{0.46\linewidth}#2\end{minipage}%
  \par\medskip
}
\providecommand{\kienthuccot}[2]{\haicotchay{#1}{#2}}
% Dự phòng nếu chưa nạp graphicx / caption (không ghi đè bản thật)
\providecommand{\resizebox}[3]{#3}
\providecommand{\captionof}[2]{\par\smallskip\noindent\textit{#2}\par}
\newenvironment{hoatdong}[1]{%
  \par\medskip\noindent\textbf{Hoạt động.} #1\par\smallskip
}{\par\medskip}
\newenvironment{emcobiet}{%
  \par\medskip\noindent\textbf{Em có biết}\par\smallskip
}{\par\medskip}
\newenvironment{emdahoc}{%
  \par\medskip\noindent\textbf{Em đã học}\par\smallskip
}{\par\medskip}
\newenvironment{emcothe}{%
  \par\medskip\noindent\textbf{Em có thể}\par\smallskip
}{\par\medskip}
\newenvironment{kienthuc}{%
  \par\medskip\noindent\textbf{Kiến thức cốt lõi}\par\smallskip
}{\par\medskip}
\newenvironment{traloi}{%
  \par\smallskip\noindent\textit{Trả lời / ghi chú của em}\par
  \vspace{1.2em}%
}{\par\medskip}
\newenvironment{phuongphap}{%
  \par\medskip\noindent\textbf{Phương pháp giải}\par\smallskip
}{\par\medskip}
\newenvironment{vidumau}{%
  \par\medskip\noindent\textbf{Bài mẫu}\par\smallskip
}{\par\medskip}
\newenvironment{dangmau}[1]{%
  \par\medskip\noindent\textbf{Dạng mẫu.} #1\par\smallskip
}{\par\medskip}
\makeatother

\makeatletter
\providecommand{\indam}[1]{\textbf{#1}}
\providecommand{\chuy}[1]{\par\smallskip\noindent\textit{#1}\par}
\providecommand{\luuy}[1]{\par\smallskip\noindent\textbf{Lưu ý.} #1\par}
\providecommand{\ghichu}[1]{\par\smallskip\noindent\textbf{Ghi chú.} #1\par}
\providecommand{\vidu}[1]{\par\smallskip\noindent\textbf{Ví dụ.} #1\par}
\providecommand{\Blythuyet}{\subsection*{Lý thuyết}}
% Hai cột: kiến thức | hình
\providecommand{\haicotchay}[2]{%
  \par\medskip\noindent
  \begin{minipage}[t]{0.52\linewidth}#1\end{minipage}\hfill
  \begin{minipage}[t]{0.46\linewidth}#2\end{minipage}%
  \par\medskip
}
\providecommand{\kienthuccot}[2]{\haicotchay{#1}{#2}}
% Dự phòng nếu chưa nạp graphicx / caption (không ghi đè bản thật)
\providecommand{\resizebox}[3]{#3}
\providecommand{\captionof}[2]{\par\smallskip\noindent\textit{#2}\par}
\newenvironment{hoatdong}[1]{%
  \par\medskip\noindent\textbf{Hoạt động.} #1\par\smallskip
}{\par\medskip}
\newenvironment{emcobiet}{%
  \par\medskip\noindent\textbf{Em có biết}\par\smallskip
}{\par\medskip}
\newenvironment{emdahoc}{%
  \par\medskip\noindent\textbf{Em đã học}\par\smallskip
}{\par\medskip}
\newenvironment{emcothe}{%
  \par\medskip\noindent\textbf{Em có thể}\par\smallskip
}{\par\medskip}
\newenvironment{kienthuc}{%
  \par\medskip\noindent\textbf{Kiến thức cốt lõi}\par\smallskip
}{\par\medskip}
\newenvironment{traloi}{%
  \par\smallskip\noindent\textit{Trả lời / ghi chú của em}\par
  \vspace{1.2em}%
}{\par\medskip}
\newenvironment{phuongphap}{%
  \par\medskip\noindent\textbf{Phương pháp giải}\par\smallskip
}{\par\medskip}
\newenvironment{vidumau}{%
  \par\medskip\noindent\textbf{Bài mẫu}\par\smallskip
}{\par\medskip}
\newenvironment{dangmau}[1]{%
  \par\medskip\noindent\textbf{Dạng mẫu.} #1\par\smallskip
}{\par\medskip}
\makeatother


\section*{L10C2 Bài 7. Đồ thị dịch chuyển - thời gian}

\subsection{Lý thuyết}

\subsubsection{Đồ thị dịch chuyển - thời gian}

\begin{kienthuc}
\textbf{Đồ thị dịch chuyển - thời gian} là đồ thị biểu diễn sự phụ thuộc của dịch chuyển $d$ của một vật vào thời gian $t$. Trục hoành biểu diễn thời gian ($t$), trục tung biểu diễn dịch chuyển ($d$).
\end{kienthuc}

\begin{hoatdong}
\textbf{Vẽ đồ thị dịch chuyển - thời gian}
\begin{enumerate}
    \item \textbf{Xác định các cặp giá trị $(t, d)$:} Từ bảng số liệu hoặc phương trình chuyển động, xác định các cặp giá trị thời gian và dịch chuyển tương ứng của vật.
    \item \textbf{Chọn tỉ lệ xích phù hợp:} Chọn tỉ lệ xích cho trục thời gian và trục dịch chuyển sao cho đồ thị dễ nhìn và thể hiện được toàn bộ quá trình chuyển động.
    \item \textbf{Vẽ các điểm trên hệ trục tọa độ:} Đánh dấu các điểm có tọa độ $(t, d)$ đã xác định lên hệ trục tọa độ.
    \item \textbf{Nối các điểm:} Nối các điểm đã đánh dấu bằng một đường cong hoặc đường thẳng. Đường này chính là đồ thị dịch chuyển - thời gian của vật.
\end{enumerate}
\end{hoatdong}

\subsubsection{Mối liên hệ giữa đồ thị dịch chuyển - thời gian và vận tốc}

\begin{kienthuc}
\textbf{Vận tốc trung bình} của một vật trong khoảng thời gian từ $t_1$ đến $t_2$ được xác định bởi công thức:
$$v_{tb} = \frac{\Delta d}{\Delta t} = \frac{d_2 - d_1}{t_2 - t_1}$$
Trên đồ thị dịch chuyển - thời gian, vận tốc trung bình trong một khoảng thời gian chính là \textbf{hệ số góc của đường nối hai điểm} tương ứng với thời điểm đầu và thời điểm cuối của khoảng thời gian đó.
\end{kienthuc}

\begin{kienthuc}
\textbf{Vận tốc tức thời} của một vật tại một thời điểm $t$ được xác định bởi công thức:
$$v = \lim_{\Delta t \to 0} \frac{\Delta d}{\Delta t}$$
Trên đồ thị dịch chuyển - thời gian, vận tốc tức thời tại một thời điểm $t$ chính là \textbf{hệ số góc của tiếp tuyến với đồ thị tại điểm tương ứng với thời điểm $t$ đó}.
\end{kienthuc}

\begin{emcobiet}
\textbf{Đồ thị dịch chuyển - thời gian cho chuyển động thẳng đều}
Đối với chuyển động thẳng đều, dịch chuyển $d$ của vật biến thiên tuyến tính theo thời gian $t$ theo phương trình $d = v \cdot t + d_0$ (nếu gốc thời gian $t=0$ và gốc tọa độ trùng nhau thì $d = v \cdot t$).
Do đó, đồ thị dịch chuyển - thời gian của chuyển động thẳng đều là một \textbf{đường thẳng}.
\begin{itemize}
    \item Nếu $v > 0$: đồ thị là đường thẳng dốc lên (hệ số góc dương).
    \item Nếu $v < 0$: đồ thị là đường thẳng dốc xuống (hệ số góc âm).
    \item Nếu $v = 0$: đồ thị là đường thẳng nằm ngang (hệ số góc bằng 0), vật đứng yên.
\end{itemize}
\end{emcobiet}

\subsubsection{Sử dụng đồ thị dịch chuyển - thời gian để mô tả chuyển động}

\begin{hoatdong}
\textbf{Phân tích đồ thị dịch chuyển - thời gian}
Từ đồ thị dịch chuyển - thời gian, ta có thể rút ra nhiều thông tin về chuyển động của vật:
\begin{enumerate}
    \item \textbf{Vị trí của vật tại các thời điểm khác nhau:} Đọc giá trị $d$ tương ứng với giá trị $t$ trên đồ thị.
    \item \textbf{Chiều chuyển động:}
    \begin{itemize}
        \item Nếu đồ thị dốc lên (hệ số góc dương), vật đang chuyển động theo chiều dương.
        \item Nếu đồ thị dốc xuống (hệ số góc âm), vật đang chuyển động theo chiều âm.
        \item Nếu đồ thị nằm ngang (hệ số góc bằng 0), vật đang đứng yên.
    \end{itemize}
    \item \textbf{Độ lớn vận tốc:} Độ lớn của hệ số góc của đồ thị (hoặc tiếp tuyến) cho biết độ lớn của vận tốc. Đường càng dốc, vận tốc càng lớn.
    \item \textbf{Sự thay đổi vận tốc:}
    \begin{itemize}
        \item Nếu đồ thị là đường thẳng, vận tốc không đổi (chuyển động thẳng đều).
        \item Nếu đồ thị là đường cong, vận tốc thay đổi (chuyển động không đều).
    \end{itemize}
\end{enumerate}
\end{hoatdong}

\begin{emcothe}
\textbf{So sánh chuyển động của hai vật bằng đồ thị}
Khi có đồ thị dịch chuyển - thời gian của hai vật trên cùng một hệ trục tọa độ, ta có thể dễ dàng so sánh:
\begin{itemize}
    \item \textbf{Thời điểm và vị trí gặp nhau:} Là giao điểm của hai đồ thị.
    \item \textbf{Vận tốc tương đối:} So sánh độ dốc của hai đồ thị.
    \item \textbf{Khoảng cách giữa hai vật:} Là hiệu các giá trị dịch chuyển tại cùng một thời điểm.
\end{itemize}
\end{emcothe}

\begin{chuy}
\textbf{Phân biệt dịch chuyển và quãng đường đi được:}
\begin{itemize}
    \item \textbf{Dịch chuyển} là đại lượng vectơ, có thể dương, âm hoặc bằng không. Trên đồ thị dịch chuyển - thời gian, nó là giá trị trên trục tung.
    \item \textbf{Quãng đường đi được} là đại lượng vô hướng, luôn không âm. Để tính quãng đường đi được từ đồ thị dịch chuyển - thời gian, ta phải tính tổng độ lớn của các dịch chuyển trong từng giai đoạn chuyển động.
\end{itemize}
\end{chuy}
